% Force the TOC to break here so Section 4 (Visual Interface) starts on a new TOC page
% with all of its subsections kept together. This keeps the TOC at exactly two pages
% with Sections 1-3 on the first page and Sections 4-6 on the second.
\addtocontents{toc}{\protect\newpage}

\section{Visual Interface}

The visual interface of Urban Scout was designed to support quick, mobile-friendly interaction for users who may be navigating the system while traveling. Since the project scenario focuses on tourists on the move, usability and readability were treated as important design concerns throughout the frontend implementation.

The application frontend was built using React with Vite, styled using Tailwind CSS, and organized into page-based views with reusable components. This stack allowed the project to support responsive layouts, fast development, and a clear separation between shared interface elements and page-specific functionality.

\subsection{Design Goals for the Interface}

The main visual design goals were:
\begin{itemize}
    \item simplicity
    \item clarity
    \item mobile responsiveness
    \item fast access to key features
    \item low navigation effort
\end{itemize}

These goals were chosen because the intended users are not expected to spend a long time learning the interface. Instead, they should be able to log in and quickly access transit, places, alerts, and itineraries with minimal friction.

\subsection{General Layout and Navigation}

The interface uses a clean and consistent layout across major pages. Important user-facing features are accessible through a persistent bottom navigation bar, which is especially useful on mobile-sized screens. This navigation style supports thumb-friendly access and reduces the need for deep menu navigation.

The main user-facing pages include: Home, Places, Transit, News, Report, Trips.

This structure was chosen so that the most important categories of information are always close at hand. Rather than nesting key features behind multiple clicks, the interface exposes them directly in the primary navigation.

In addition to regular user pages, the system includes separate admin pages for tasks such as report verification and news management. This helps preserve role-based separation while keeping each interface focused on the needs of its specific user type.

\begin{figure}[H]
    \centering
    \IfFileExists{figures/screenshots/01-login.png}{%
        \includegraphics[width=0.7\textwidth,keepaspectratio]{figures/screenshots/01-login.png}%
    }{%
        \fbox{\parbox{0.7\textwidth}{\centering\vspace{2cm}\textbf{Placeholder for 01-login.png}\\\vspace{0.3cm}Login page with email field, User/Admin toggle, and Log In button.\vspace{2cm}}}%
    }
    \caption{Login screen with the User and Admin role toggle.}
    \label{fig:login}
\end{figure}

\subsection{Home Page}

The home page acts as the central landing page after login. It introduces the user to the application and displays region-related information in a clear card-based format. Because regions are an important organizing concept in the database, surfacing them on the home page helps users understand how the system is structured.

The home page also includes summary information and a clear entry point into the rest of the application. Its purpose is not only decorative but functional: it provides orientation and gives the user a quick sense of what is available.

\begin{figure}[H]
    \centering
    \IfFileExists{figures/screenshots/02-home.png}{%
        \includegraphics[width=0.7\textwidth,keepaspectratio]{figures/screenshots/02-home.png}%
    }{%
        \fbox{\parbox{0.7\textwidth}{\centering\vspace{2cm}\textbf{Placeholder for 02-home.png}\\\vspace{0.3cm}Hero banner, three region cards with Subscribe or Subscribed buttons, and a stats row.\vspace{2cm}}}%
    }
    \caption{Home page showing region cards with subscribe and unsubscribe controls.}
    \label{fig:home}
\end{figure}

\subsection{Places Interface}

The places interface was designed to help users browse destinations in an intuitive way. Users can filter places by region and category, such as landmarks, eateries, or transit stations. This supports practical tourism use cases such as: finding attractions in a region, locating food options nearby, identifying station-based destinations.

Cards are used to display places because they allow important information such as names, categories, and addresses to be scanned quickly. This format is more user-friendly than presenting the same information as dense text or raw table data.

\begin{figure}[H]
    \centering
    \IfFileExists{figures/screenshots/03-places.png}{%
        \includegraphics[width=0.7\textwidth,keepaspectratio]{figures/screenshots/03-places.png}%
    }{%
        \fbox{\parbox{0.7\textwidth}{\centering\vspace{2cm}\textbf{Placeholder for 03-places.png}\\\vspace{0.3cm}Type filter pills (Landmarks, Eateries, Transit, All), region dropdown, and a grid of place cards.\vspace{2cm}}}%
    }
    \caption{Places browse view with type filter pills and the region dropdown.}
    \label{fig:places}
\end{figure}

\begin{figure}[H]
    \centering
    \IfFileExists{figures/screenshots/04-place-detail.png}{%
        \includegraphics[width=0.7\textwidth,keepaspectratio]{figures/screenshots/04-place-detail.png}%
    }{%
        \fbox{\parbox{0.7\textwidth}{\centering\vspace{2cm}\textbf{Placeholder for 04-place-detail.png}\\\vspace{0.3cm}Hero image, address card, Leaflet map with marker, nearby transit list, and Add to Itinerary controls.\vspace{2cm}}}%
    }
    \caption{Place detail page with the embedded map and nearby transit list.}
    \label{fig:placedetail}
\end{figure}

\subsection{Transit Interface}

The transit page presents route and station-related information in a structured format. Users can select a station and view arrivals, departures, serving routes, and nearby places. Because transit data can become visually overwhelming if poorly presented, the page organizes information into grouped sections with clear labels.
\nopagebreak[4]

This page is one of the strongest examples of the project's integrated design. It combines multiple layers of information that would normally require separate applications or manual effort from the user. The interface presents these relationships in a way that is easier to understand than raw backend data.

\begin{figure}[H]
    \centering
    \IfFileExists{figures/screenshots/05-transit.png}{%
        \includegraphics[width=0.7\textwidth,keepaspectratio]{figures/screenshots/05-transit.png}%
    }{%
        \fbox{\parbox{0.7\textwidth}{\centering\vspace{2cm}\textbf{Placeholder for 05-transit.png}\\\vspace{0.3cm}Routes list, station dropdown selected, arrivals and departures tables, and nearby places cards.\vspace{2cm}}}%
    }
    \caption{Transit page showing arrivals, departures, and nearby places for a selected station.}
    \label{fig:transit}
\end{figure}

\subsection{News and Safety Alerts Interface}

The news page was designed to support both awareness and filtering. Users can browse regional news posts and view associated safety alerts. The interface presents each news post in a card-based layout with important metadata such as severity, date, region, and linked alert information.

This page is especially important from a usability perspective because it adds a local awareness layer to the tourism experience. A tourist may not only want to know what places exist, but also whether there are closures, crowds, hazards, or events that affect their plans. The visual grouping of posts and alerts helps users process this information more efficiently.

\begin{figure}[H]
    \centering
    \IfFileExists{figures/screenshots/06-news.png}{%
        \includegraphics[width=0.7\textwidth,keepaspectratio]{figures/screenshots/06-news.png}%
    }{%
        \fbox{\parbox{0.7\textwidth}{\centering\vspace{2cm}\textbf{Placeholder for 06-news.png}\\\vspace{0.3cm}Region filter pills, news cards with severity badges, and one card showing the red safety alert box.\vspace{2cm}}}%
    }
    \caption{News feed with a region filter and a safety alert card.}
    \label{fig:news}
\end{figure}

\subsection{Incident Report Interface}

The report submission page uses a form-based layout with structured input fields for region, category, severity, and description. This page was designed to keep user input straightforward and focused. Because many users may only submit a report once or occasionally, the form avoids unnecessary complexity and aims for clarity.

In a complete use scenario, the usability of this form matters because it affects whether users are willing to contribute useful information to the system. The design therefore emphasizes simple field selection, recognizable labels, and a direct submission flow.

\begin{figure}[H]
    \centering
    \IfFileExists{figures/screenshots/07-incident.png}{%
        \includegraphics[width=0.7\textwidth,keepaspectratio]{figures/screenshots/07-incident.png}%
    }{%
        \fbox{\parbox{0.7\textwidth}{\centering\vspace{2cm}\textbf{Placeholder for 07-incident.png}\\\vspace{0.3cm}Two-column layout with the New Report form on one side (region dropdown, category dropdown, severity pills, description) and My Reports on the other.\vspace{2cm}}}%
    }
    \caption{Incident report form alongside the user's previously submitted reports.}
    \label{fig:incident}
\end{figure}

\subsection{Itinerary Interface}

The itinerary pages allow users to view, create, and manage travel plans. The design is centered around a list of itineraries and detail views for specific trips. Since itineraries are planning-oriented features, the interface aims to be structured and easy to scan rather than visually dense.

The itinerary feature helps demonstrate that the interface supports not only information browsing, but also user-driven organization of travel activity.

\begin{figure}[H]
    \centering
    \IfFileExists{figures/screenshots/08-itinerary.png}{%
        \includegraphics[width=0.7\textwidth,keepaspectratio]{figures/screenshots/08-itinerary.png}%
    }{%
        \fbox{\parbox{0.7\textwidth}{\centering\vspace{2cm}\textbf{Placeholder for 08-itinerary.png}\\\vspace{0.3cm}List of trips with create form expanded, and an itinerary detail showing item cards with Edit and Remove buttons.\vspace{2cm}}}%
    }
    \caption{Itinerary list view and an open trip detail page.}
    \label{fig:itinerary}
\end{figure}

\subsection{Admin Interface}

The admin interface was designed separately from the tourist-facing side of the application. It includes pages for dashboard-style management, incident report verification, news management, place management, and transit route management. These pages focus more on clarity and efficiency than on tourism-style browsing.

This separation improves usability because administrators and tourists have different goals. Administrators need direct access to moderation and content tasks, while regular users need location-based exploration and planning tools.

\begin{figure}[H]
    \centering
    \IfFileExists{figures/screenshots/09-admin-dashboard.png}{%
        \includegraphics[width=0.7\textwidth,keepaspectratio]{figures/screenshots/09-admin-dashboard.png}%
    }{%
        \fbox{\parbox{0.7\textwidth}{\centering\vspace{2cm}\textbf{Placeholder for 09-admin-dashboard.png}\\\vspace{0.3cm}Five stat cards (Pending Reports, News Posts, Places, Routes, Total Reports) and a recent pending reports preview.\vspace{2cm}}}%
    }
    \caption{Admin dashboard with stat cards for the assigned region.}
    \label{fig:admindash}
\end{figure}

\begin{figure}[H]
    \centering
    \IfFileExists{figures/screenshots/10-admin-news.png}{%
        \includegraphics[width=0.7\textwidth,keepaspectratio]{figures/screenshots/10-admin-news.png}%
    }{%
        \fbox{\parbox{0.7\textwidth}{\centering\vspace{2cm}\textbf{Placeholder for 10-admin-news.png}\\\vspace{0.3cm}Form with title, severity pills, body textarea, Include Safety Alert checked, alert type dropdown visible, and Publish button.\vspace{2cm}}}%
    }
    \caption{Admin Manage News page with the new post form open and a safety alert attached.}
    \label{fig:adminnews}
\end{figure}

\subsection{Interface Friendliness and Usability}

Several qualities contribute to the friendliness and usability of the interface:

\boldheading{Consistency} the application uses consistent page structure, spacing, and navigation patterns across major views. This helps users move between pages without relearning the interface.

\boldheading{Readability} information is grouped into cards, lists, and labeled sections rather than being displayed as raw data. This makes the application much easier to use for non-technical users.

\boldheading{Mobile-First Design} the layout and navigation were designed with smaller screens in mind. Because the intended users are tourists, this is an especially important design strength.

\boldheading{Low Learning Curve} the interface avoids excessive feature depth and instead focuses on the main flows users are likely to need. This makes the app easier to understand quickly.

\boldheading{Connection to Real Tasks} each page supports a recognizable real-world action, such as browsing attractions, checking transit, reading local updates, reporting incidents, or planning a trip. This makes the interface feel purposeful rather than artificially constructed.

\subsection{Minimum Keyboard Input}

A core usability priority for Urban Scout was to minimize how much the user has to type. Since the target audience is tourists using the app on a phone while moving through a city, the interface was built around selection-based controls like dropdowns, pill toggles, native date pickers, and tap-to-add buttons rather than free-text fields wherever possible.

\begin{description}
    \item[Region selection] Dropdown menus on the Home, Places, News, Incidents, and Trips pages let the user pick a region from a known list instead of typing it.
    \item[Place type filter] Pill toggle buttons on the Places page (Landmarks, Eateries, Transit, All) let the user switch categories with a single tap.
    \item[Severity] Pill-shaped Low, Medium, and High buttons on the Incident Report form and on the admin Manage News form remove the need to type a severity level.
    \item[Category] A dropdown on the Incident Report form lets the user pick from preset incident categories rather than typing one.
    \item[Station] A dropdown on the Transit page lets the user select a station from the list of stations in the chosen region.
    \item[Alert type] A dropdown on the admin Manage News form is used when a safety alert is included in a news post, so the alert type is selected rather than typed.
    \item[Trip dates] Native HTML date pickers on the Trips form are used for start and end dates, so the user never has to type a date string manually.
    \item[Itinerary picker] On the Place Detail page, the user picks an existing itinerary from a dropdown and presses a single Add button to attach the place, which avoids any typing during itinerary building.
\end{description}

The only fields that genuinely require typing are the incident description, the news title, and the news body, since these are free-text content fields where a fixed selection list is not possible.

\subsection{Summary}

The visual interface of Urban Scout was designed to support a unified, usable, and mobile-friendly experience. It connects the database-backed features of the system to a set of pages that are practical for both tourists and administrators. The interface emphasizes clarity, direct access to major functions, and readable presentation of linked information. As a result, it supports the overall project goal of reducing fragmentation in travel-related information access.
